\chapter*{Introduction Générale}
\addcontentsline{toc}{chapter}{Introduction Générale}
\markboth{Introduction Générale}{Introduction Générale}

Le cloud computing est devenu le socle de la recherche scientifique et de
l'enseignement supérieur modernes. La capacité à provisionner des ressources
informatiques à la demande --- machines virtuelles, conteneurs et environnements
cluster --- supprime les obstacles traditionnels liés aux cycles d'acquisition de
matériel et permet aux chercheurs d'itérer plus rapidement. En Tunisie, le Ministère
de l'Enseignement Supérieur et de la Recherche Scientifique (MESRS) opère une
infrastructure cloud nationale Huawei Cloud Stack (HCS) à l'adresse
\textit{mesrscloud.rnu.tn}, mettant des ressources cloud privées à la disposition
des institutions académiques à travers le pays. Cependant, l'accès à ces ressources
est resté un processus manuel et technique : les utilisateurs doivent soumettre des
demandes aux administrateurs, attendre le provisionnement et recevoir leurs identifiants
sans aucune visibilité sur le processus et sans autonomie de libre-service.

\textbf{Mission de stage.} Ce projet de fin d'études a été réalisé chez Huawei
Technologies Tunisie avec la mission suivante : concevoir et développer une plateforme
cloud en libre-service --- \textit{VM Marketplace} --- permettant aux utilisateurs
authentifiés de l'environnement HCS MESRS de configurer, payer, provisionner,
superviser et gérer de manière autonome des machines virtuelles et des clusters
Kubernetes via une interface web intuitive, sans intervention d'un administrateur pour
les opérations courantes.

\textbf{Problématique.} La problématique centrale de ce projet est la suivante :
\textit{Comment permettre aux utilisateurs académiques de l'infrastructure cloud HCS
du MESRS de provisionner et gérer de façon autonome des machines virtuelles et des
clusters Kubernetes via une plateforme conviviale intégrant paiement, recommandation
IA, suivi en temps réel et supervision automatisée ?}

\textbf{Objectifs.} Pour répondre à cette problématique, cinq objectifs techniques
ont été définis :
\begin{enumerate}
  \item Implémenter une application web full-stack sécurisée avec authentification
        utilisateur, assistant de provisionnement de VM en plusieurs étapes, intégration
        du paiement Stripe et infrastructure as code Terraform sur HCS.
  \item Développer un moteur de recommandation IA qui traduit des descriptions
        textuelles de charge de travail en spécifications de VM optimales à l'aide
        d'un LLM hébergé localement.
  \item Ajouter un terminal SSH dans le navigateur et une installation automatisée
        de Netdata pour chaque VM provisionnée.
  \item Concevoir et implémenter un pipeline de provisionnement de clusters Kubernetes
        qui initialise des clusters multi-nœuds sur HCS avec kubeadm, en incluant une
        solution réseau permettant aux nœuds workers d'accéder à Internet malgré les
        contraintes de quota EIP.
  \item Valider l'ensemble de la plateforme sur l'infrastructure HCS de production.
\end{enumerate}

\textbf{Structure du rapport.} Ce rapport est organisé en quatre chapitres suivis
d'une conclusion générale.

\begin{itemize}
  \item \textbf{Chapitre 1 --- Contexte général} présente l'organisme d'accueil,
        situe le projet dans l'écosystème cloud académique du MESRS, étudie les
        solutions existantes et justifie l'approche proposée et la méthodologie.

  \item \textbf{Chapitre 2 --- Analyse et Architecture} identifie les acteurs et
        les besoins, établit le backlog produit et détaille l'architecture à trois
        services et les choix technologiques.

  \item \textbf{Chapitre 3 --- Conception et Réalisation} documente l'implémentation
        à travers trois sprints itératifs : provisionnement de VM et authentification,
        recommandation IA et supervision, et gestion des clusters Kubernetes.

  \item \textbf{Chapitre 4 --- Tests et Validation} présente les résultats des tests
        et évalue la plateforme au regard des objectifs initiaux.
\end{itemize}
